\documentclass[aspectratio=169]{beamer}
% Shared preamble for ml/ slide decks.
% Each lecture .tex starts with:
%   \documentclass[aspectratio=169]{beamer}
%   % Shared preamble for ml/ slide decks.
% Each lecture .tex starts with:
%   \documentclass[aspectratio=169]{beamer}
%   % Shared preamble for ml/ slide decks.
% Each lecture .tex starts with:
%   \documentclass[aspectratio=169]{beamer}
%   \input{../preamble}
% Compile from inside the chapter folder so \input{../preamble} resolves.

\usetheme{default}
\usecolortheme{dove}

\setbeamertemplate{navigation symbols}{}
\setbeamertemplate{footline}{%
  \hfill{\large\insertframenumber\,/\,\inserttotalframenumber}\hspace{0.8em}\vspace{0.5em}%
}

% Project palette (existing, kept for backwards compatibility with deferred/*)
\definecolor{popblue}{RGB}{52, 101, 164}
\definecolor{sampred}{RGB}{204, 0, 0}
\definecolor{paramgreen}{RGB}{0, 140, 70}
\definecolor{warnred}{RGB}{180, 40, 40}
\definecolor{orange1}{RGB}{220, 120, 0}
\definecolor{violet1}{RGB}{120, 50, 160}
\definecolor{lightbg}{RGB}{245, 245, 250}

% Armenian flag palette (use these for any 3-color highlight; per CLAUDE.md).
% "Nice versions": slightly de-saturated from the official spec for slide use,
% so they don't burn out under projector colors. Originals kept as *_pure.
\definecolor{armred}{RGB}{200, 30, 40}      % nice red (official #D90012 muted)
\definecolor{armblue}{RGB}{30, 70, 160}     % nice blue (official #0033A0 muted)
\definecolor{armorange}{RGB}{230, 160, 30}  % nice orange (official #F2A800 muted)
\definecolor{armred_pure}{HTML}{D90012}     % official Armenian flag red
\definecolor{armblue_pure}{HTML}{0033A0}    % official Armenian flag blue
\definecolor{armorange_pure}{HTML}{F2A800}  % official Armenian flag orange

\setbeamercolor{frametitle}{fg=popblue}
\setbeamercolor{title}{fg=popblue}

\usepackage{pgfplots}
\usepackage{tikz}
\usetikzlibrary{shapes, arrows.meta, positioning, calc, decorations.pathreplacing, patterns, matrix}
\pgfplotsset{compat=1.18}
\usepackage{amsmath, amssymb, amsfonts}
\usepackage{bm}
\usepackage{dsfont}
\usepackage{array}
\usepackage{pifont}
\usepackage{colortbl}
\usepackage{fontenc}

% Callout boxes. Purely additive - existing decks are unaffected.
% Usage: \begin{tcolorbox}[colback=armblue!8, colframe=armblue, boxrule=0.6pt, arc=2pt,
%                          title=\textbf{Key idea}] ... \end{tcolorbox}
% Palette convention (ml/SLIDE_STYLE.md): armblue=key idea, armred=trap/warning,
% paramgreen=takeaway or "Next", armorange=watch-out.
\usepackage[most]{tcolorbox}

% Cleaner table rules (\toprule, \midrule, \bottomrule). Also purely additive.
\usepackage{booktabs}

% \cancel{...} for struck-through terms in derivations (added 2026-08-11 for the
% importance-sampling frame in ch11_rl, where the dynamics cancel). Purely additive.
\usepackage{cancel}

% Minimal ML notation macros (subset of slds-lmu latex-math/basic-ml.tex).
% Kept here so we don't have to inline-expand every \thetav, \xv, etc.
% when porting upstream slides.
\newcommand{\R}{\mathds{R}}
\newcommand{\N}{\mathds{N}}
\newcommand{\E}{\mathds{E}}
\newcommand{\Xspace}{\mathcal{X}}
\newcommand{\Yspace}{\mathcal{Y}}
\newcommand{\Hspace}{\mathcal{H}}
\newcommand{\D}{\mathcal{D}}
\newcommand{\Dtrain}{\mathcal{D}_{\text{train}}}
\newcommand{\Dtest}{\mathcal{D}_{\text{test}}}
\newcommand{\xv}{\mathbf{x}}
\newcommand{\yv}{\mathbf{y}}
\newcommand{\Xmat}{\mathbf{X}}
\newcommand{\id}{\mathbf{I}}
\newcommand{\thetav}{\bm{\theta}}
\newcommand{\thetah}{\hat{\bm{\theta}}}
\newcommand{\lamv}{\bm{\lambda}}
\newcommand{\fx}{f(\xv)}
\newcommand{\fxt}{f(\xv \,|\, \bm{\theta})}
\newcommand{\fxi}{f\!\left(\xv^{(i)}\right)}
\newcommand{\fh}{\hat{f}}
\newcommand{\fxh}{\hat{f}(\xv)}
\newcommand{\yh}{\hat{y}}
\newcommand{\yih}{\hat{y}^{(i)}}
% \xi clashes with the Greek letter \xi; write \xv^{(i)} inline instead.
\newcommand{\yi}[1][i]{y^{(#1)}}
\newcommand{\xyi}[1][i]{\left(\xv^{(#1)},\, y^{(#1)}\right)}
\newcommand{\thx}{\bm{\theta}^\top \xv}
\newcommand{\sumin}{\sum_{i=1}^{n}}
\newcommand{\sumjp}{\sum_{j=1}^{p}}
\newcommand{\Lxy}{L\!\left(y, \fx\right)}
\newcommand{\Lxyi}{L\!\left(\yi, \fxi\right)}
\newcommand{\risk}{\mathcal{R}}
\newcommand{\riske}{\mathcal{R}_{\text{emp}}}
\newcommand{\risket}{\mathcal{R}_{\text{emp}}(\bm{\theta})}
\newcommand{\riskt}{\mathcal{R}(\bm{\theta})}
\newcommand{\riskr}{\mathcal{R}_{\text{reg}}}
\newcommand{\argmin}{\mathop{\mathrm{arg\,min}}}
\newcommand{\argmax}{\mathop{\mathrm{arg\,max}}}
\newcommand{\eps}{\varepsilon}

% Compile from inside the chapter folder so % Shared preamble for ml/ slide decks.
% Each lecture .tex starts with:
%   \documentclass[aspectratio=169]{beamer}
%   \input{../preamble}
% Compile from inside the chapter folder so \input{../preamble} resolves.

\usetheme{default}
\usecolortheme{dove}

\setbeamertemplate{navigation symbols}{}
\setbeamertemplate{footline}{%
  \hfill{\large\insertframenumber\,/\,\inserttotalframenumber}\hspace{0.8em}\vspace{0.5em}%
}

% Project palette (existing, kept for backwards compatibility with deferred/*)
\definecolor{popblue}{RGB}{52, 101, 164}
\definecolor{sampred}{RGB}{204, 0, 0}
\definecolor{paramgreen}{RGB}{0, 140, 70}
\definecolor{warnred}{RGB}{180, 40, 40}
\definecolor{orange1}{RGB}{220, 120, 0}
\definecolor{violet1}{RGB}{120, 50, 160}
\definecolor{lightbg}{RGB}{245, 245, 250}

% Armenian flag palette (use these for any 3-color highlight; per CLAUDE.md).
% "Nice versions": slightly de-saturated from the official spec for slide use,
% so they don't burn out under projector colors. Originals kept as *_pure.
\definecolor{armred}{RGB}{200, 30, 40}      % nice red (official #D90012 muted)
\definecolor{armblue}{RGB}{30, 70, 160}     % nice blue (official #0033A0 muted)
\definecolor{armorange}{RGB}{230, 160, 30}  % nice orange (official #F2A800 muted)
\definecolor{armred_pure}{HTML}{D90012}     % official Armenian flag red
\definecolor{armblue_pure}{HTML}{0033A0}    % official Armenian flag blue
\definecolor{armorange_pure}{HTML}{F2A800}  % official Armenian flag orange

\setbeamercolor{frametitle}{fg=popblue}
\setbeamercolor{title}{fg=popblue}

\usepackage{pgfplots}
\usepackage{tikz}
\usetikzlibrary{shapes, arrows.meta, positioning, calc, decorations.pathreplacing, patterns, matrix}
\pgfplotsset{compat=1.18}
\usepackage{amsmath, amssymb, amsfonts}
\usepackage{bm}
\usepackage{dsfont}
\usepackage{array}
\usepackage{pifont}
\usepackage{colortbl}
\usepackage{fontenc}

% Callout boxes. Purely additive - existing decks are unaffected.
% Usage: \begin{tcolorbox}[colback=armblue!8, colframe=armblue, boxrule=0.6pt, arc=2pt,
%                          title=\textbf{Key idea}] ... \end{tcolorbox}
% Palette convention (ml/SLIDE_STYLE.md): armblue=key idea, armred=trap/warning,
% paramgreen=takeaway or "Next", armorange=watch-out.
\usepackage[most]{tcolorbox}

% Cleaner table rules (\toprule, \midrule, \bottomrule). Also purely additive.
\usepackage{booktabs}

% \cancel{...} for struck-through terms in derivations (added 2026-08-11 for the
% importance-sampling frame in ch11_rl, where the dynamics cancel). Purely additive.
\usepackage{cancel}

% Minimal ML notation macros (subset of slds-lmu latex-math/basic-ml.tex).
% Kept here so we don't have to inline-expand every \thetav, \xv, etc.
% when porting upstream slides.
\newcommand{\R}{\mathds{R}}
\newcommand{\N}{\mathds{N}}
\newcommand{\E}{\mathds{E}}
\newcommand{\Xspace}{\mathcal{X}}
\newcommand{\Yspace}{\mathcal{Y}}
\newcommand{\Hspace}{\mathcal{H}}
\newcommand{\D}{\mathcal{D}}
\newcommand{\Dtrain}{\mathcal{D}_{\text{train}}}
\newcommand{\Dtest}{\mathcal{D}_{\text{test}}}
\newcommand{\xv}{\mathbf{x}}
\newcommand{\yv}{\mathbf{y}}
\newcommand{\Xmat}{\mathbf{X}}
\newcommand{\id}{\mathbf{I}}
\newcommand{\thetav}{\bm{\theta}}
\newcommand{\thetah}{\hat{\bm{\theta}}}
\newcommand{\lamv}{\bm{\lambda}}
\newcommand{\fx}{f(\xv)}
\newcommand{\fxt}{f(\xv \,|\, \bm{\theta})}
\newcommand{\fxi}{f\!\left(\xv^{(i)}\right)}
\newcommand{\fh}{\hat{f}}
\newcommand{\fxh}{\hat{f}(\xv)}
\newcommand{\yh}{\hat{y}}
\newcommand{\yih}{\hat{y}^{(i)}}
% \xi clashes with the Greek letter \xi; write \xv^{(i)} inline instead.
\newcommand{\yi}[1][i]{y^{(#1)}}
\newcommand{\xyi}[1][i]{\left(\xv^{(#1)},\, y^{(#1)}\right)}
\newcommand{\thx}{\bm{\theta}^\top \xv}
\newcommand{\sumin}{\sum_{i=1}^{n}}
\newcommand{\sumjp}{\sum_{j=1}^{p}}
\newcommand{\Lxy}{L\!\left(y, \fx\right)}
\newcommand{\Lxyi}{L\!\left(\yi, \fxi\right)}
\newcommand{\risk}{\mathcal{R}}
\newcommand{\riske}{\mathcal{R}_{\text{emp}}}
\newcommand{\risket}{\mathcal{R}_{\text{emp}}(\bm{\theta})}
\newcommand{\riskt}{\mathcal{R}(\bm{\theta})}
\newcommand{\riskr}{\mathcal{R}_{\text{reg}}}
\newcommand{\argmin}{\mathop{\mathrm{arg\,min}}}
\newcommand{\argmax}{\mathop{\mathrm{arg\,max}}}
\newcommand{\eps}{\varepsilon}
 resolves.

\usetheme{default}
\usecolortheme{dove}

\setbeamertemplate{navigation symbols}{}
\setbeamertemplate{footline}{%
  \hfill{\large\insertframenumber\,/\,\inserttotalframenumber}\hspace{0.8em}\vspace{0.5em}%
}

% Project palette (existing, kept for backwards compatibility with deferred/*)
\definecolor{popblue}{RGB}{52, 101, 164}
\definecolor{sampred}{RGB}{204, 0, 0}
\definecolor{paramgreen}{RGB}{0, 140, 70}
\definecolor{warnred}{RGB}{180, 40, 40}
\definecolor{orange1}{RGB}{220, 120, 0}
\definecolor{violet1}{RGB}{120, 50, 160}
\definecolor{lightbg}{RGB}{245, 245, 250}

% Armenian flag palette (use these for any 3-color highlight; per CLAUDE.md).
% "Nice versions": slightly de-saturated from the official spec for slide use,
% so they don't burn out under projector colors. Originals kept as *_pure.
\definecolor{armred}{RGB}{200, 30, 40}      % nice red (official #D90012 muted)
\definecolor{armblue}{RGB}{30, 70, 160}     % nice blue (official #0033A0 muted)
\definecolor{armorange}{RGB}{230, 160, 30}  % nice orange (official #F2A800 muted)
\definecolor{armred_pure}{HTML}{D90012}     % official Armenian flag red
\definecolor{armblue_pure}{HTML}{0033A0}    % official Armenian flag blue
\definecolor{armorange_pure}{HTML}{F2A800}  % official Armenian flag orange

\setbeamercolor{frametitle}{fg=popblue}
\setbeamercolor{title}{fg=popblue}

\usepackage{pgfplots}
\usepackage{tikz}
\usetikzlibrary{shapes, arrows.meta, positioning, calc, decorations.pathreplacing, patterns, matrix}
\pgfplotsset{compat=1.18}
\usepackage{amsmath, amssymb, amsfonts}
\usepackage{bm}
\usepackage{dsfont}
\usepackage{array}
\usepackage{pifont}
\usepackage{colortbl}
\usepackage{fontenc}

% Callout boxes. Purely additive - existing decks are unaffected.
% Usage: \begin{tcolorbox}[colback=armblue!8, colframe=armblue, boxrule=0.6pt, arc=2pt,
%                          title=\textbf{Key idea}] ... \end{tcolorbox}
% Palette convention (ml/SLIDE_STYLE.md): armblue=key idea, armred=trap/warning,
% paramgreen=takeaway or "Next", armorange=watch-out.
\usepackage[most]{tcolorbox}

% Cleaner table rules (\toprule, \midrule, \bottomrule). Also purely additive.
\usepackage{booktabs}

% \cancel{...} for struck-through terms in derivations (added 2026-08-11 for the
% importance-sampling frame in ch11_rl, where the dynamics cancel). Purely additive.
\usepackage{cancel}

% Minimal ML notation macros (subset of slds-lmu latex-math/basic-ml.tex).
% Kept here so we don't have to inline-expand every \thetav, \xv, etc.
% when porting upstream slides.
\newcommand{\R}{\mathds{R}}
\newcommand{\N}{\mathds{N}}
\newcommand{\E}{\mathds{E}}
\newcommand{\Xspace}{\mathcal{X}}
\newcommand{\Yspace}{\mathcal{Y}}
\newcommand{\Hspace}{\mathcal{H}}
\newcommand{\D}{\mathcal{D}}
\newcommand{\Dtrain}{\mathcal{D}_{\text{train}}}
\newcommand{\Dtest}{\mathcal{D}_{\text{test}}}
\newcommand{\xv}{\mathbf{x}}
\newcommand{\yv}{\mathbf{y}}
\newcommand{\Xmat}{\mathbf{X}}
\newcommand{\id}{\mathbf{I}}
\newcommand{\thetav}{\bm{\theta}}
\newcommand{\thetah}{\hat{\bm{\theta}}}
\newcommand{\lamv}{\bm{\lambda}}
\newcommand{\fx}{f(\xv)}
\newcommand{\fxt}{f(\xv \,|\, \bm{\theta})}
\newcommand{\fxi}{f\!\left(\xv^{(i)}\right)}
\newcommand{\fh}{\hat{f}}
\newcommand{\fxh}{\hat{f}(\xv)}
\newcommand{\yh}{\hat{y}}
\newcommand{\yih}{\hat{y}^{(i)}}
% \xi clashes with the Greek letter \xi; write \xv^{(i)} inline instead.
\newcommand{\yi}[1][i]{y^{(#1)}}
\newcommand{\xyi}[1][i]{\left(\xv^{(#1)},\, y^{(#1)}\right)}
\newcommand{\thx}{\bm{\theta}^\top \xv}
\newcommand{\sumin}{\sum_{i=1}^{n}}
\newcommand{\sumjp}{\sum_{j=1}^{p}}
\newcommand{\Lxy}{L\!\left(y, \fx\right)}
\newcommand{\Lxyi}{L\!\left(\yi, \fxi\right)}
\newcommand{\risk}{\mathcal{R}}
\newcommand{\riske}{\mathcal{R}_{\text{emp}}}
\newcommand{\risket}{\mathcal{R}_{\text{emp}}(\bm{\theta})}
\newcommand{\riskt}{\mathcal{R}(\bm{\theta})}
\newcommand{\riskr}{\mathcal{R}_{\text{reg}}}
\newcommand{\argmin}{\mathop{\mathrm{arg\,min}}}
\newcommand{\argmax}{\mathop{\mathrm{arg\,max}}}
\newcommand{\eps}{\varepsilon}

% Compile from inside the chapter folder so % Shared preamble for ml/ slide decks.
% Each lecture .tex starts with:
%   \documentclass[aspectratio=169]{beamer}
%   % Shared preamble for ml/ slide decks.
% Each lecture .tex starts with:
%   \documentclass[aspectratio=169]{beamer}
%   \input{../preamble}
% Compile from inside the chapter folder so \input{../preamble} resolves.

\usetheme{default}
\usecolortheme{dove}

\setbeamertemplate{navigation symbols}{}
\setbeamertemplate{footline}{%
  \hfill{\large\insertframenumber\,/\,\inserttotalframenumber}\hspace{0.8em}\vspace{0.5em}%
}

% Project palette (existing, kept for backwards compatibility with deferred/*)
\definecolor{popblue}{RGB}{52, 101, 164}
\definecolor{sampred}{RGB}{204, 0, 0}
\definecolor{paramgreen}{RGB}{0, 140, 70}
\definecolor{warnred}{RGB}{180, 40, 40}
\definecolor{orange1}{RGB}{220, 120, 0}
\definecolor{violet1}{RGB}{120, 50, 160}
\definecolor{lightbg}{RGB}{245, 245, 250}

% Armenian flag palette (use these for any 3-color highlight; per CLAUDE.md).
% "Nice versions": slightly de-saturated from the official spec for slide use,
% so they don't burn out under projector colors. Originals kept as *_pure.
\definecolor{armred}{RGB}{200, 30, 40}      % nice red (official #D90012 muted)
\definecolor{armblue}{RGB}{30, 70, 160}     % nice blue (official #0033A0 muted)
\definecolor{armorange}{RGB}{230, 160, 30}  % nice orange (official #F2A800 muted)
\definecolor{armred_pure}{HTML}{D90012}     % official Armenian flag red
\definecolor{armblue_pure}{HTML}{0033A0}    % official Armenian flag blue
\definecolor{armorange_pure}{HTML}{F2A800}  % official Armenian flag orange

\setbeamercolor{frametitle}{fg=popblue}
\setbeamercolor{title}{fg=popblue}

\usepackage{pgfplots}
\usepackage{tikz}
\usetikzlibrary{shapes, arrows.meta, positioning, calc, decorations.pathreplacing, patterns, matrix}
\pgfplotsset{compat=1.18}
\usepackage{amsmath, amssymb, amsfonts}
\usepackage{bm}
\usepackage{dsfont}
\usepackage{array}
\usepackage{pifont}
\usepackage{colortbl}
\usepackage{fontenc}

% Callout boxes. Purely additive - existing decks are unaffected.
% Usage: \begin{tcolorbox}[colback=armblue!8, colframe=armblue, boxrule=0.6pt, arc=2pt,
%                          title=\textbf{Key idea}] ... \end{tcolorbox}
% Palette convention (ml/SLIDE_STYLE.md): armblue=key idea, armred=trap/warning,
% paramgreen=takeaway or "Next", armorange=watch-out.
\usepackage[most]{tcolorbox}

% Cleaner table rules (\toprule, \midrule, \bottomrule). Also purely additive.
\usepackage{booktabs}

% \cancel{...} for struck-through terms in derivations (added 2026-08-11 for the
% importance-sampling frame in ch11_rl, where the dynamics cancel). Purely additive.
\usepackage{cancel}

% Minimal ML notation macros (subset of slds-lmu latex-math/basic-ml.tex).
% Kept here so we don't have to inline-expand every \thetav, \xv, etc.
% when porting upstream slides.
\newcommand{\R}{\mathds{R}}
\newcommand{\N}{\mathds{N}}
\newcommand{\E}{\mathds{E}}
\newcommand{\Xspace}{\mathcal{X}}
\newcommand{\Yspace}{\mathcal{Y}}
\newcommand{\Hspace}{\mathcal{H}}
\newcommand{\D}{\mathcal{D}}
\newcommand{\Dtrain}{\mathcal{D}_{\text{train}}}
\newcommand{\Dtest}{\mathcal{D}_{\text{test}}}
\newcommand{\xv}{\mathbf{x}}
\newcommand{\yv}{\mathbf{y}}
\newcommand{\Xmat}{\mathbf{X}}
\newcommand{\id}{\mathbf{I}}
\newcommand{\thetav}{\bm{\theta}}
\newcommand{\thetah}{\hat{\bm{\theta}}}
\newcommand{\lamv}{\bm{\lambda}}
\newcommand{\fx}{f(\xv)}
\newcommand{\fxt}{f(\xv \,|\, \bm{\theta})}
\newcommand{\fxi}{f\!\left(\xv^{(i)}\right)}
\newcommand{\fh}{\hat{f}}
\newcommand{\fxh}{\hat{f}(\xv)}
\newcommand{\yh}{\hat{y}}
\newcommand{\yih}{\hat{y}^{(i)}}
% \xi clashes with the Greek letter \xi; write \xv^{(i)} inline instead.
\newcommand{\yi}[1][i]{y^{(#1)}}
\newcommand{\xyi}[1][i]{\left(\xv^{(#1)},\, y^{(#1)}\right)}
\newcommand{\thx}{\bm{\theta}^\top \xv}
\newcommand{\sumin}{\sum_{i=1}^{n}}
\newcommand{\sumjp}{\sum_{j=1}^{p}}
\newcommand{\Lxy}{L\!\left(y, \fx\right)}
\newcommand{\Lxyi}{L\!\left(\yi, \fxi\right)}
\newcommand{\risk}{\mathcal{R}}
\newcommand{\riske}{\mathcal{R}_{\text{emp}}}
\newcommand{\risket}{\mathcal{R}_{\text{emp}}(\bm{\theta})}
\newcommand{\riskt}{\mathcal{R}(\bm{\theta})}
\newcommand{\riskr}{\mathcal{R}_{\text{reg}}}
\newcommand{\argmin}{\mathop{\mathrm{arg\,min}}}
\newcommand{\argmax}{\mathop{\mathrm{arg\,max}}}
\newcommand{\eps}{\varepsilon}

% Compile from inside the chapter folder so % Shared preamble for ml/ slide decks.
% Each lecture .tex starts with:
%   \documentclass[aspectratio=169]{beamer}
%   \input{../preamble}
% Compile from inside the chapter folder so \input{../preamble} resolves.

\usetheme{default}
\usecolortheme{dove}

\setbeamertemplate{navigation symbols}{}
\setbeamertemplate{footline}{%
  \hfill{\large\insertframenumber\,/\,\inserttotalframenumber}\hspace{0.8em}\vspace{0.5em}%
}

% Project palette (existing, kept for backwards compatibility with deferred/*)
\definecolor{popblue}{RGB}{52, 101, 164}
\definecolor{sampred}{RGB}{204, 0, 0}
\definecolor{paramgreen}{RGB}{0, 140, 70}
\definecolor{warnred}{RGB}{180, 40, 40}
\definecolor{orange1}{RGB}{220, 120, 0}
\definecolor{violet1}{RGB}{120, 50, 160}
\definecolor{lightbg}{RGB}{245, 245, 250}

% Armenian flag palette (use these for any 3-color highlight; per CLAUDE.md).
% "Nice versions": slightly de-saturated from the official spec for slide use,
% so they don't burn out under projector colors. Originals kept as *_pure.
\definecolor{armred}{RGB}{200, 30, 40}      % nice red (official #D90012 muted)
\definecolor{armblue}{RGB}{30, 70, 160}     % nice blue (official #0033A0 muted)
\definecolor{armorange}{RGB}{230, 160, 30}  % nice orange (official #F2A800 muted)
\definecolor{armred_pure}{HTML}{D90012}     % official Armenian flag red
\definecolor{armblue_pure}{HTML}{0033A0}    % official Armenian flag blue
\definecolor{armorange_pure}{HTML}{F2A800}  % official Armenian flag orange

\setbeamercolor{frametitle}{fg=popblue}
\setbeamercolor{title}{fg=popblue}

\usepackage{pgfplots}
\usepackage{tikz}
\usetikzlibrary{shapes, arrows.meta, positioning, calc, decorations.pathreplacing, patterns, matrix}
\pgfplotsset{compat=1.18}
\usepackage{amsmath, amssymb, amsfonts}
\usepackage{bm}
\usepackage{dsfont}
\usepackage{array}
\usepackage{pifont}
\usepackage{colortbl}
\usepackage{fontenc}

% Callout boxes. Purely additive - existing decks are unaffected.
% Usage: \begin{tcolorbox}[colback=armblue!8, colframe=armblue, boxrule=0.6pt, arc=2pt,
%                          title=\textbf{Key idea}] ... \end{tcolorbox}
% Palette convention (ml/SLIDE_STYLE.md): armblue=key idea, armred=trap/warning,
% paramgreen=takeaway or "Next", armorange=watch-out.
\usepackage[most]{tcolorbox}

% Cleaner table rules (\toprule, \midrule, \bottomrule). Also purely additive.
\usepackage{booktabs}

% \cancel{...} for struck-through terms in derivations (added 2026-08-11 for the
% importance-sampling frame in ch11_rl, where the dynamics cancel). Purely additive.
\usepackage{cancel}

% Minimal ML notation macros (subset of slds-lmu latex-math/basic-ml.tex).
% Kept here so we don't have to inline-expand every \thetav, \xv, etc.
% when porting upstream slides.
\newcommand{\R}{\mathds{R}}
\newcommand{\N}{\mathds{N}}
\newcommand{\E}{\mathds{E}}
\newcommand{\Xspace}{\mathcal{X}}
\newcommand{\Yspace}{\mathcal{Y}}
\newcommand{\Hspace}{\mathcal{H}}
\newcommand{\D}{\mathcal{D}}
\newcommand{\Dtrain}{\mathcal{D}_{\text{train}}}
\newcommand{\Dtest}{\mathcal{D}_{\text{test}}}
\newcommand{\xv}{\mathbf{x}}
\newcommand{\yv}{\mathbf{y}}
\newcommand{\Xmat}{\mathbf{X}}
\newcommand{\id}{\mathbf{I}}
\newcommand{\thetav}{\bm{\theta}}
\newcommand{\thetah}{\hat{\bm{\theta}}}
\newcommand{\lamv}{\bm{\lambda}}
\newcommand{\fx}{f(\xv)}
\newcommand{\fxt}{f(\xv \,|\, \bm{\theta})}
\newcommand{\fxi}{f\!\left(\xv^{(i)}\right)}
\newcommand{\fh}{\hat{f}}
\newcommand{\fxh}{\hat{f}(\xv)}
\newcommand{\yh}{\hat{y}}
\newcommand{\yih}{\hat{y}^{(i)}}
% \xi clashes with the Greek letter \xi; write \xv^{(i)} inline instead.
\newcommand{\yi}[1][i]{y^{(#1)}}
\newcommand{\xyi}[1][i]{\left(\xv^{(#1)},\, y^{(#1)}\right)}
\newcommand{\thx}{\bm{\theta}^\top \xv}
\newcommand{\sumin}{\sum_{i=1}^{n}}
\newcommand{\sumjp}{\sum_{j=1}^{p}}
\newcommand{\Lxy}{L\!\left(y, \fx\right)}
\newcommand{\Lxyi}{L\!\left(\yi, \fxi\right)}
\newcommand{\risk}{\mathcal{R}}
\newcommand{\riske}{\mathcal{R}_{\text{emp}}}
\newcommand{\risket}{\mathcal{R}_{\text{emp}}(\bm{\theta})}
\newcommand{\riskt}{\mathcal{R}(\bm{\theta})}
\newcommand{\riskr}{\mathcal{R}_{\text{reg}}}
\newcommand{\argmin}{\mathop{\mathrm{arg\,min}}}
\newcommand{\argmax}{\mathop{\mathrm{arg\,max}}}
\newcommand{\eps}{\varepsilon}
 resolves.

\usetheme{default}
\usecolortheme{dove}

\setbeamertemplate{navigation symbols}{}
\setbeamertemplate{footline}{%
  \hfill{\large\insertframenumber\,/\,\inserttotalframenumber}\hspace{0.8em}\vspace{0.5em}%
}

% Project palette (existing, kept for backwards compatibility with deferred/*)
\definecolor{popblue}{RGB}{52, 101, 164}
\definecolor{sampred}{RGB}{204, 0, 0}
\definecolor{paramgreen}{RGB}{0, 140, 70}
\definecolor{warnred}{RGB}{180, 40, 40}
\definecolor{orange1}{RGB}{220, 120, 0}
\definecolor{violet1}{RGB}{120, 50, 160}
\definecolor{lightbg}{RGB}{245, 245, 250}

% Armenian flag palette (use these for any 3-color highlight; per CLAUDE.md).
% "Nice versions": slightly de-saturated from the official spec for slide use,
% so they don't burn out under projector colors. Originals kept as *_pure.
\definecolor{armred}{RGB}{200, 30, 40}      % nice red (official #D90012 muted)
\definecolor{armblue}{RGB}{30, 70, 160}     % nice blue (official #0033A0 muted)
\definecolor{armorange}{RGB}{230, 160, 30}  % nice orange (official #F2A800 muted)
\definecolor{armred_pure}{HTML}{D90012}     % official Armenian flag red
\definecolor{armblue_pure}{HTML}{0033A0}    % official Armenian flag blue
\definecolor{armorange_pure}{HTML}{F2A800}  % official Armenian flag orange

\setbeamercolor{frametitle}{fg=popblue}
\setbeamercolor{title}{fg=popblue}

\usepackage{pgfplots}
\usepackage{tikz}
\usetikzlibrary{shapes, arrows.meta, positioning, calc, decorations.pathreplacing, patterns, matrix}
\pgfplotsset{compat=1.18}
\usepackage{amsmath, amssymb, amsfonts}
\usepackage{bm}
\usepackage{dsfont}
\usepackage{array}
\usepackage{pifont}
\usepackage{colortbl}
\usepackage{fontenc}

% Callout boxes. Purely additive - existing decks are unaffected.
% Usage: \begin{tcolorbox}[colback=armblue!8, colframe=armblue, boxrule=0.6pt, arc=2pt,
%                          title=\textbf{Key idea}] ... \end{tcolorbox}
% Palette convention (ml/SLIDE_STYLE.md): armblue=key idea, armred=trap/warning,
% paramgreen=takeaway or "Next", armorange=watch-out.
\usepackage[most]{tcolorbox}

% Cleaner table rules (\toprule, \midrule, \bottomrule). Also purely additive.
\usepackage{booktabs}

% \cancel{...} for struck-through terms in derivations (added 2026-08-11 for the
% importance-sampling frame in ch11_rl, where the dynamics cancel). Purely additive.
\usepackage{cancel}

% Minimal ML notation macros (subset of slds-lmu latex-math/basic-ml.tex).
% Kept here so we don't have to inline-expand every \thetav, \xv, etc.
% when porting upstream slides.
\newcommand{\R}{\mathds{R}}
\newcommand{\N}{\mathds{N}}
\newcommand{\E}{\mathds{E}}
\newcommand{\Xspace}{\mathcal{X}}
\newcommand{\Yspace}{\mathcal{Y}}
\newcommand{\Hspace}{\mathcal{H}}
\newcommand{\D}{\mathcal{D}}
\newcommand{\Dtrain}{\mathcal{D}_{\text{train}}}
\newcommand{\Dtest}{\mathcal{D}_{\text{test}}}
\newcommand{\xv}{\mathbf{x}}
\newcommand{\yv}{\mathbf{y}}
\newcommand{\Xmat}{\mathbf{X}}
\newcommand{\id}{\mathbf{I}}
\newcommand{\thetav}{\bm{\theta}}
\newcommand{\thetah}{\hat{\bm{\theta}}}
\newcommand{\lamv}{\bm{\lambda}}
\newcommand{\fx}{f(\xv)}
\newcommand{\fxt}{f(\xv \,|\, \bm{\theta})}
\newcommand{\fxi}{f\!\left(\xv^{(i)}\right)}
\newcommand{\fh}{\hat{f}}
\newcommand{\fxh}{\hat{f}(\xv)}
\newcommand{\yh}{\hat{y}}
\newcommand{\yih}{\hat{y}^{(i)}}
% \xi clashes with the Greek letter \xi; write \xv^{(i)} inline instead.
\newcommand{\yi}[1][i]{y^{(#1)}}
\newcommand{\xyi}[1][i]{\left(\xv^{(#1)},\, y^{(#1)}\right)}
\newcommand{\thx}{\bm{\theta}^\top \xv}
\newcommand{\sumin}{\sum_{i=1}^{n}}
\newcommand{\sumjp}{\sum_{j=1}^{p}}
\newcommand{\Lxy}{L\!\left(y, \fx\right)}
\newcommand{\Lxyi}{L\!\left(\yi, \fxi\right)}
\newcommand{\risk}{\mathcal{R}}
\newcommand{\riske}{\mathcal{R}_{\text{emp}}}
\newcommand{\risket}{\mathcal{R}_{\text{emp}}(\bm{\theta})}
\newcommand{\riskt}{\mathcal{R}(\bm{\theta})}
\newcommand{\riskr}{\mathcal{R}_{\text{reg}}}
\newcommand{\argmin}{\mathop{\mathrm{arg\,min}}}
\newcommand{\argmax}{\mathop{\mathrm{arg\,max}}}
\newcommand{\eps}{\varepsilon}
 resolves.

\usetheme{default}
\usecolortheme{dove}

\setbeamertemplate{navigation symbols}{}
\setbeamertemplate{footline}{%
  \hfill{\large\insertframenumber\,/\,\inserttotalframenumber}\hspace{0.8em}\vspace{0.5em}%
}

% Project palette (existing, kept for backwards compatibility with deferred/*)
\definecolor{popblue}{RGB}{52, 101, 164}
\definecolor{sampred}{RGB}{204, 0, 0}
\definecolor{paramgreen}{RGB}{0, 140, 70}
\definecolor{warnred}{RGB}{180, 40, 40}
\definecolor{orange1}{RGB}{220, 120, 0}
\definecolor{violet1}{RGB}{120, 50, 160}
\definecolor{lightbg}{RGB}{245, 245, 250}

% Armenian flag palette (use these for any 3-color highlight; per CLAUDE.md).
% "Nice versions": slightly de-saturated from the official spec for slide use,
% so they don't burn out under projector colors. Originals kept as *_pure.
\definecolor{armred}{RGB}{200, 30, 40}      % nice red (official #D90012 muted)
\definecolor{armblue}{RGB}{30, 70, 160}     % nice blue (official #0033A0 muted)
\definecolor{armorange}{RGB}{230, 160, 30}  % nice orange (official #F2A800 muted)
\definecolor{armred_pure}{HTML}{D90012}     % official Armenian flag red
\definecolor{armblue_pure}{HTML}{0033A0}    % official Armenian flag blue
\definecolor{armorange_pure}{HTML}{F2A800}  % official Armenian flag orange

\setbeamercolor{frametitle}{fg=popblue}
\setbeamercolor{title}{fg=popblue}

\usepackage{pgfplots}
\usepackage{tikz}
\usetikzlibrary{shapes, arrows.meta, positioning, calc, decorations.pathreplacing, patterns, matrix}
\pgfplotsset{compat=1.18}
\usepackage{amsmath, amssymb, amsfonts}
\usepackage{bm}
\usepackage{dsfont}
\usepackage{array}
\usepackage{pifont}
\usepackage{colortbl}
\usepackage{fontenc}

% Callout boxes. Purely additive - existing decks are unaffected.
% Usage: \begin{tcolorbox}[colback=armblue!8, colframe=armblue, boxrule=0.6pt, arc=2pt,
%                          title=\textbf{Key idea}] ... \end{tcolorbox}
% Palette convention (ml/SLIDE_STYLE.md): armblue=key idea, armred=trap/warning,
% paramgreen=takeaway or "Next", armorange=watch-out.
\usepackage[most]{tcolorbox}

% Cleaner table rules (\toprule, \midrule, \bottomrule). Also purely additive.
\usepackage{booktabs}

% \cancel{...} for struck-through terms in derivations (added 2026-08-11 for the
% importance-sampling frame in ch11_rl, where the dynamics cancel). Purely additive.
\usepackage{cancel}

% Minimal ML notation macros (subset of slds-lmu latex-math/basic-ml.tex).
% Kept here so we don't have to inline-expand every \thetav, \xv, etc.
% when porting upstream slides.
\newcommand{\R}{\mathds{R}}
\newcommand{\N}{\mathds{N}}
\newcommand{\E}{\mathds{E}}
\newcommand{\Xspace}{\mathcal{X}}
\newcommand{\Yspace}{\mathcal{Y}}
\newcommand{\Hspace}{\mathcal{H}}
\newcommand{\D}{\mathcal{D}}
\newcommand{\Dtrain}{\mathcal{D}_{\text{train}}}
\newcommand{\Dtest}{\mathcal{D}_{\text{test}}}
\newcommand{\xv}{\mathbf{x}}
\newcommand{\yv}{\mathbf{y}}
\newcommand{\Xmat}{\mathbf{X}}
\newcommand{\id}{\mathbf{I}}
\newcommand{\thetav}{\bm{\theta}}
\newcommand{\thetah}{\hat{\bm{\theta}}}
\newcommand{\lamv}{\bm{\lambda}}
\newcommand{\fx}{f(\xv)}
\newcommand{\fxt}{f(\xv \,|\, \bm{\theta})}
\newcommand{\fxi}{f\!\left(\xv^{(i)}\right)}
\newcommand{\fh}{\hat{f}}
\newcommand{\fxh}{\hat{f}(\xv)}
\newcommand{\yh}{\hat{y}}
\newcommand{\yih}{\hat{y}^{(i)}}
% \xi clashes with the Greek letter \xi; write \xv^{(i)} inline instead.
\newcommand{\yi}[1][i]{y^{(#1)}}
\newcommand{\xyi}[1][i]{\left(\xv^{(#1)},\, y^{(#1)}\right)}
\newcommand{\thx}{\bm{\theta}^\top \xv}
\newcommand{\sumin}{\sum_{i=1}^{n}}
\newcommand{\sumjp}{\sum_{j=1}^{p}}
\newcommand{\Lxy}{L\!\left(y, \fx\right)}
\newcommand{\Lxyi}{L\!\left(\yi, \fxi\right)}
\newcommand{\risk}{\mathcal{R}}
\newcommand{\riske}{\mathcal{R}_{\text{emp}}}
\newcommand{\risket}{\mathcal{R}_{\text{emp}}(\bm{\theta})}
\newcommand{\riskt}{\mathcal{R}(\bm{\theta})}
\newcommand{\riskr}{\mathcal{R}_{\text{reg}}}
\newcommand{\argmin}{\mathop{\mathrm{arg\,min}}}
\newcommand{\argmax}{\mathop{\mathrm{arg\,max}}}
\newcommand{\eps}{\varepsilon}

\newcolumntype{L}[1]{>{\raggedright\arraybackslash}p{#1}}

\title{Lecture 37b: TabPFN After the Original Paper}
\subtitle{Scale, new data types, and the cost of a forward pass}
\date{}

\begin{document}

\begin{frame}\titlepage\end{frame}

\begin{frame}{The row limit in L37 was a snapshot}
  \small
  The Nature 2025 paper made its central quality claim on tables with up to
  \textbf{10,000 rows}. It did not establish a permanent limit for prior-data fitted networks.
  \vskip 12pt
  \begin{center}
  \begin{tabular}{lccc}
    \textbf{Release} & \textbf{v2} & \textbf{2.5} & \textbf{3 / 3.5}\\[3pt]
    Design scale reported & 10k & 50k & up to 1m\\
    Core story & small tables & wider tables & cached, chunked inference
  \end{tabular}
  \end{center}
  \vskip 10pt
  \fcolorbox{armorange}{armorange!12}{\parbox{0.96\linewidth}{\small
    A design range, a hardware capacity, and a benchmark win are three different claims.
    Each number on the following slides keeps its paper and comparison attached.}}
  \vskip 5pt
  {\scriptsize Sources: Hollmann et al. (2025); Prior Labs (2025, 2026a, 2026b).}
\end{frame}

\begin{frame}{Outline}\tableofcontents\end{frame}

\section{The release line}
\begin{frame}[plain]
  \vfill\centering
  {\Large\color{popblue}\textbf{The release line}}\\[9pt]
  {\small What was actually changed after the first useful model?}
  \vfill
\end{frame}

\begin{frame}{Five checkpoints, three distinct questions}
  \small
  \begin{tabular}{L{0.16\linewidth}L{0.29\linewidth}L{0.43\linewidth}}
    \textbf{Paper} & \textbf{Problem targeted} & \textbf{Main development}\\[3pt]\hline
    v1, 2022 & tiny numerical classification & proof that prior-data fitting works\\[3pt]
    v2, 2025 & practical small tables & cell representations and two-way attention\\[3pt]
    2.5, 2025 & more rows and features & stronger predictor and output distillation\\[3pt]
    3, 2026 & large contexts and new tasks & cache, row chunking, many-class decoder\\[3pt]
    3.5, 2026 & messy and shifted tables & revised encodings, prior, and model width
  \end{tabular}
  \vskip 10pt
  \textbf{Ask separately:} How large a context fits? Which input types work? What wins on a specified split?
  \vskip 5pt
  {\scriptsize Sources: Hollmann et al. (2022, 2025); Prior Labs (2025, 2026a, 2026b).}
\end{frame}

\section{TabPFN-2.5}
\begin{frame}[plain]
  \vfill\centering
  {\Large\color{popblue}\textbf{TabPFN-2.5}}\\[9pt]
  {\small More of the table fits in context.}
  \vfill
\end{frame}

\begin{frame}{Fifty thousand rows and two thousand features}
  \small
  The TabPFN-2.5 report describes a model \textbf{built for up to 50,000 data points and 2,000
  features}, a 20-fold increase in data cells over v2.
  \vskip 10pt
  \begin{columns}[T]
    \column{0.49\textwidth}
    \fcolorbox{armblue}{armblue!8}{\parbox{0.94\linewidth}{\small
      \textbf{Designed range}\par
      50k rows, 2k features. These are model design targets, not a guarantee on every
      device or dataset.}}
    \column{0.49\textwidth}
    \fcolorbox{paramgreen}{paramgreen!8}{\parbox{0.94\linewidth}{\small
      \textbf{Benchmark reach}\par
      The report also evaluates datasets with up to 100k rows. That is an empirical
      result beyond the design target.}}
  \end{columns}
  \vskip 10pt
  It reports strong default performance against XGBoost and compares with an AutoGluon
  ensemble tuned for four hours. Those are \textbf{benchmark-specific} comparisons.
  \vskip 4pt
  {\scriptsize Source: Prior Labs, \emph{TabPFN-2.5} (2025), arXiv:2511.08667.}
\end{frame}

\begin{frame}{A fast fit can still have expensive repeated inference}
  \small
  TabPFN conditions on labeled rows when predicting. If an application scores many new rows,
  repeated context processing can dominate latency.
  The report distills the predictor into a multilayer perceptron (MLP) or tree ensemble.
  \vskip 9pt
  \begin{center}
  \begin{tikzpicture}[>=stealth, node distance=1.35cm, every node/.style={font=\small}]
    \node[draw=popblue,fill=armblue!8,rounded corners,minimum width=2.5cm,minimum height=.9cm] (p) {TabPFN-2.5};
    \node[draw=armorange,fill=armorange!12,rounded corners,right of=p,xshift=2.3cm,minimum width=2.4cm,minimum height=.9cm] (d) {distillation};
    \node[draw=paramgreen,fill=paramgreen!8,rounded corners,right of=d,xshift=2.3cm,minimum width=2.6cm,minimum height=.9cm] (s) {MLP or trees};
    \draw[->,thick] (p) -- (d);
    \draw[->,thick] (d) -- (s);
  \end{tikzpicture}
  \end{center}
  \vskip 7pt
  The report's \textbf{TabPFN-as-MLP/TreeEns} output engine trades some fidelity to
  the teacher for much lower serving latency. Distillation creates a new predictor;
  it is not the original forward pass made magically cheap.
  \vskip 4pt
  {\scriptsize Source: Prior Labs (2025), arXiv:2511.08667. Schematic, not measured latency.}
\end{frame}

\section{TabPFN-3}
\begin{frame}[plain]
  \vfill\centering
  {\Large\color{popblue}\textbf{TabPFN-3}}\\[9pt]
  {\small Reuse the work on the context; process rows in chunks.}
  \vfill
\end{frame}

\begin{frame}{One million rows needs a different inference path}
  \small
  The TabPFN-3 report describes \textbf{row chunking} and a reduced key-value cache
  that let a single H100 handle up to \textbf{one million training rows}.
  \vskip 10pt
  \begin{center}
  \begin{tikzpicture}[>=stealth, every node/.style={font=\small}]
    \node[draw=popblue,fill=armblue!8,rounded corners,minimum width=3.3cm,minimum height=.8cm] (a) at (0,0) {labeled table};
    \node[draw=armorange,fill=armorange!12,rounded corners,minimum width=3.3cm,minimum height=.8cm] (b) at (4.2,0) {row chunks};
    \node[draw=paramgreen,fill=paramgreen!8,rounded corners,minimum width=3.3cm,minimum height=.8cm] (c) at (8.4,0) {cached context};
    \node[draw=popblue,fill=armblue!8,rounded corners,minimum width=3.3cm,minimum height=.8cm] (d) at (8.4,-1.5) {many test queries};
    \draw[->,thick] (a) -- (b); \draw[->,thick] (b) -- (c); \draw[->,thick] (c) -- (d);
  \end{tikzpicture}
  \end{center}
  \vskip 7pt
  Cache the training context once, then reuse it while scoring test rows. This addresses
  repeated inference cost. It does \emph{not} make a million-row table cheap on arbitrary hardware.
  \vskip 4pt
  {\scriptsize Source: Prior Labs, \emph{TabPFN-3} (2026), arXiv:2605.13986. Schematic.}
\end{frame}

\begin{frame}{The task frontier widened, but through different routes}
  \small
  \begin{tabular}{L{0.29\linewidth}L{0.62\linewidth}}
    \textbf{Setting} & \textbf{What the TabPFN-3 report actually describes}\\[3pt]\hline
    Many-class classification & A new decoder for larger class counts.\\[4pt]
    Time series & A specialized checkpoint, TabPFN-TS-3; second on fev-bench in the report.\\[4pt]
    Relational data & A TabPFN-based approach evaluated on RelBenchV1.\\[4pt]
    Text in tables & TabPFN-3-Plus on TabSTAR, rather than a blanket claim for every checkpoint.
  \end{tabular}
  \vskip 11pt
  \fcolorbox{armorange}{armorange!12}{\parbox{0.96\linewidth}{\small
    A task-specific harness or Plus service is a different product path from giving raw text
    to the base numerical/categorical model.}}
  \vskip 4pt
  {\scriptsize Source: Prior Labs (2026), arXiv:2605.13986.}
\end{frame}

\begin{frame}{Thinking mode spends compute at prediction time}
  \small
  L37's basic prediction is one fixed forward pass. TabPFN-3 also introduces a
  \textbf{Thinking} mode that spends more computation at inference to improve accuracy.
  \vskip 12pt
  \begin{center}
  \begin{tikzpicture}[>=stealth, every node/.style={font=\small}]
    \node[draw=popblue,fill=armblue!8,rounded corners,minimum width=3.7cm,minimum height=1cm] (a) at (0,0) {default forward pass};
    \node[draw=armorange,fill=armorange!12,rounded corners,minimum width=3.7cm,minimum height=1cm] (b) at (5.3,0) {more test-time compute};
    \draw[->,very thick] (a) -- (b);
  \end{tikzpicture}
  \end{center}
  \vskip 12pt
  The report's headline Thinking results describe an \textbf{API offering}. Keep its
  accuracy and access claims separate from the local base checkpoint.
  \vskip 4pt
  {\scriptsize Source: Prior Labs (2026), arXiv:2605.13986.}
\end{frame}

\section{TabPFN-3.5}
\begin{frame}[plain]
  \vfill\centering
  {\Large\color{popblue}\textbf{TabPFN-3.5}}\\[9pt]
  {\small Test beyond clean, independent rows.}
  \vfill
\end{frame}

\begin{frame}{The new target is the awkward real table}
  \small
  The 3.5 report explicitly targets failures revealed by broader benchmarks:
  \vskip 8pt
  \begin{itemize}\itemsep 7pt
    \item \textbf{Grouped or temporal splits}: training and test rows are not interchangeable.
    \item \textbf{Strings and high-cardinality categories}: simple numeric encodings lose information.
    \item \textbf{Wide tables and mixed modalities}: many features, text and images complicate prediction.
  \end{itemize}
  \vskip 10pt
  Its reported improvements do not erase the need to evaluate on the \emph{deployment split}.
  A random split can make a method look better than a later-time or unseen-group test.
  \vskip 5pt
  {\scriptsize Source: Prior Labs, \emph{TabPFN-3.5} (2026), arXiv:2609.17895, introduction.}
\end{frame}

\begin{frame}{What changed inside 3.5?}
  \small
  The report names three changes behind the base model's gains:
  \vskip 10pt
  \begin{enumerate}\itemsep 8pt
    \item \textbf{Revised per-cell encodings}: the representation of an individual value changes.
    \item \textbf{Wider model}: model width doubles while the cache size stays unchanged.
    \item \textbf{Revised synthetic prior}: pretraining emphasizes wide tables,
      high-cardinality features, and grouped data.
  \end{enumerate}
  \vskip 10pt
  These changes preserve L37's core idea: train on tasks sampled from a prior, then
  condition on the user's table. The prior itself moves as the target data move.
  \vskip 4pt
  {\scriptsize Source: Prior Labs (2026), arXiv:2609.17895, introduction and Section 3.}
\end{frame}

\begin{frame}{Core, Fast, Plus and Thinking are different claims}
  \small
  \begin{tabular}{L{0.23\linewidth}L{0.34\linewidth}L{0.34\linewidth}}
    \textbf{Variant} & \textbf{What it is for} & \textbf{Careful reading}\\[3pt]\hline
    3.5 core & Standard classification and regression & One jointly trained checkpoint.\\[4pt]
    3.5-Fast & Lower latency & Reported up to $3\times$ faster than 3, with most gains retained.\\[4pt]
    3.5-Plus & Richer text/date handling & Additional multimodal and proprietary optimizations.\\[4pt]
    3.5-Thinking & More inference computation & Different accuracy--latency point, not the base forward pass.
  \end{tabular}
  \vskip 8pt
  \textbf{Input check:} numeric/categorical/missing values are core tabular inputs;
  raw strings, text, images and dates require the appropriate pipeline or variant.
  \vskip 4pt
  {\scriptsize Source: Prior Labs (2026), arXiv:2609.17895, abstract and introduction.}
\end{frame}

\begin{frame}{How to read the seven-benchmark headline}
  \small
  The 3.5 report states \textbf{rank 1 on seven benchmark families}. The table's
  leading TabPFN member is \emph{not the same checkpoint in every row}:
  \vskip 8pt
  \begin{tabular}{L{0.29\linewidth}L{0.63\linewidth}}
    Standard TabArena and several other suites & 3.5-Thinking leads in the paper's summary.\\[5pt]
    RelArena-$\alpha$ & A relational variant, TabPFN-Rel (3.5), leads.\\[5pt]
    TALENT and ScoringBench & The base 3.5 model leads in the summary table.
  \end{tabular}
  \vskip 11pt
  The paper also reports strong results on BeyondArena, STRABLE and MulTaBench.
  These benchmarks probe different data types and split rules; a single rank cannot
  predict which model wins on a new table.
  \vskip 4pt
  {\scriptsize Source: Prior Labs (2026), arXiv:2609.17895, Table 1.}
\end{frame}

\section{A follow-up for learning}
\begin{frame}[plain]
  \vfill\centering
  {\Large\color{popblue}\textbf{A follow-up for learning}}\\[9pt]
  {\small A smaller implementation helps us see the mechanism.}
  \vfill
\end{frame}

\begin{frame}{nanoTabPFN makes the training loop inspectable}
  \small
  Pfefferle, Hog, Purucker and Hutter (2025) reimplement the \textbf{v2 architecture}
  with a simpler training loop and pre-generated synthetic tasks.
  \vskip 10pt
  \begin{columns}[T]
    \column{0.49\textwidth}
    \fcolorbox{armblue}{armblue!8}{\parbox{0.94\linewidth}{\small
      \textbf{Use it to study}\par
      How tasks are sampled, how examples enter context, and how one network learns
      to predict labels on new tasks.}}
    \column{0.49\textwidth}
    \fcolorbox{armorange}{armorange!12}{\parbox{0.94\linewidth}{\small
      \textbf{Keep its claim bounded}\par
      One minute of training on a GPU reaches traditional-baseline performance
      \emph{in a restricted small-data setting}.}}
  \end{columns}
  \vskip 10pt
  This is a teaching and research implementation, not a replacement for the
  latest checkpoint or its reported benchmark performance.
  \vskip 4pt
  {\scriptsize Source: Pfefferle et al. (2025), arXiv:2511.03634.}
\end{frame}

\section*{Wrap-up}
\begin{frame}{Recap}
  \small
  \begin{enumerate}\itemsep 7pt
    \item \textbf{2.5} expanded the designed table size and offered distillation for serving.
    \item \textbf{3} used caching and row chunking to reach million-row contexts on specified hardware.
    \item \textbf{3.5} targeted less tidy data and revised its encodings and synthetic prior.
    \item \textbf{Variants matter}: base, Fast, Plus, Thinking and task-specific harnesses
      have different capabilities and access paths.
    \item \textbf{nanoTabPFN} exposes the training recipe at educational scale.
  \end{enumerate}
  \vskip 7pt
  \fcolorbox{paramgreen}{paramgreen!8}{\parbox{0.96\linewidth}{\small
    \textbf{Next: L37c.} If a model can condition on a modified table in one pass,
    what does that make possible for explanation?}}
\end{frame}

\end{document}

% Provenance:
%   Sources: Hollmann et al. 2022 arXiv:2207.01848; Hollmann et al. 2025
%   Nature DOI:10.1038/s41586-024-08328-6; Prior Labs 2025 arXiv:2511.08667;
%   Prior Labs 2026 arXiv:2605.13986 and arXiv:2609.17895;
%   Pfefferle et al. 2025 arXiv:2511.03634. PDFs are in papers/.
%   Diagrams are schematic TikZ illustrations, not measured results.
%   Forward pointer: L37c_tabpfn_explanations.tex.
